\Needspace{35\baselineskip}\SourceLeaf{84}

\begin{center}
I.
\end{center}

\noindent\hspace*{\parindent}Since for $u = 0$ we have $x = t$, \textsc{Taylor}'s theorem gives

\[
\varphi x = \varphi t + u \frac{d\varphi x}{du} + \frac{uu}{1.2}\frac{d^2\varphi x}{du^2}
+ \frac{u^3}{1.2.3}\frac{d^3\varphi x}{du^3} + \text{etc.},
\]

\Needspace{23\baselineskip}\noindent where the values of the quantities $\dfrac{d\varphi x}{du}$, $\dfrac{d^2\varphi x}{du^2}$, etc.\ are to be taken at
$u = 0$. In general, however,

\[
\frac{d\varphi x}{du} = X\frac{d\varphi x}{dt},
\]

\[
\frac{d^2\varphi x}{du^2} = \frac{dX^2 \frac{d\varphi x}{dt}}{dt},
\]

\[
\frac{d^3\varphi x}{du^3} = \frac{d^2X^3 \frac{d\varphi x}{dt}}{dt^2},
\]

\[
\cdot\quad\cdot\quad\cdot\quad\cdot\quad\cdot\quad\cdot\quad\cdot\quad\cdot\quad\cdot,
\]

\noindent as I shall prove in a moment. It is also easy to see that for $u = 0$, i.\,e.\
$x = t$, the following expressions change as indicated,

\[
X\frac{d\varphi x}{dt},\quad \frac{dX^2 \frac{d\varphi x}{dt}}{dt},\quad
\text{etc.\ into}\quad X\frac{d\varphi x}{dx},\quad \frac{dX^2 \frac{d\varphi x}{dx}}{dx},\quad \text{etc.}
\]

\noindent with $x = t$ then to be substituted in the resulting expressions. This gives

\[
\varphi x = \varphi t + uX\frac{d\varphi x}{dx} + \frac{uu}{1.2}\frac{dX^2 \frac{d\varphi x}{dx}}{dx}
+ \frac{u^3}{1.2.3}\frac{d^2X^3 \frac{d\varphi x}{dx}}{dx^2} + \cdot\,\cdot\,\cdot,
\]

\noindent where $x = t$ is to be substituted, and this is Mr. \textsc{de la Grange}'s formula.

\medskip

\Needspace{10\baselineskip}\begin{center}
II.
\end{center}

\noindent\hspace*{\parindent}I prove the transformations used here as follows:

\noindent\hspace*{\parindent}I differentiate the given equation $x = t + uX$ with respect to $t$ and to $u$,
which gives

\[
\frac{dx}{dt} = 1 + u\frac{dX}{dt} = 1 + u\frac{dX}{dx}\frac{dx}{dt},
\]

\[
\frac{dx}{du} = X + u\frac{dX}{du} = X + u\frac{dX}{dx}\frac{dx}{du}.
\]
